\section*{Chapter \ref{ch:b2:metric} --- Topology of Metric Spaces}

\begin{solution}{exo:b2:metric:1}
$\delta$: symmetry and separation are clear; triangle inequality:
$\min(1, u + v) \leq \min(1,u) + \min(1,v)$ for $u, v \geq 0$ (if
either min is $1$, the right side is $\geq 1$; else it is $u + v$).
$d$: it is the pullback of $\abs{\cdot}$ by the injective $\arctan$:
the three axioms transfer.

Convergence: for $\delta$, $\delta(x_n, x) \to 0 \iff \abs{x_n - x}
\to 0$ (for small values the two distances agree): same convergent
sequences as usual. For $d$: $d(x_n, x) \to 0 \iff \arctan x_n \to
\arctan x \iff x_n \to x$ (continuity and strict monotonicity of
$\arctan$ and of its inverse on the relevant ranges): again the
usual convergence.

$(\R, d)$ is \emph{not} complete: $x_n = n$ satisfies $d(x_p, x_q)
= \abs{\arctan p - \arctan q} \to 0$ (both tend to $\frac\pi2$):
Cauchy; but $(x_n)$ does not converge for $d$ (its $d$-limit would
be an ordinary limit). Completeness is a property of the
\emph{distance}, not just of the convergent sequences.
\end{solution}

\begin{solution}{exo:b2:metric:2}
Convergent $\Rightarrow$ Cauchy: $d(x_p, x_q) \leq d(x_p, \ell) +
d(\ell, x_q)$. Bounded: beyond $N$, $d(x_n, \ell) \leq 1$; the
finitely many first terms are within some radius too.

Cauchy $+$ convergent subsequence $x_{\varphi(n)} \to \ell$: given
$\varepsilon$, for large $n$, $d(x_n, \ell) \leq d(x_n,
x_{\varphi(n)}) + d(x_{\varphi(n)}, \ell) \leq 2\varepsilon$ (the
first term by Cauchy, since $\varphi(n) \geq n$).

Compact $\Rightarrow$ complete: a Cauchy sequence has a convergent
subsequence (compactness), hence converges.
\end{solution}

\begin{solution}{exo:b2:metric:3}
$d_\infty(f, g) = \sup_{\intcc{0}{1}} \abs{x - x^2} = \frac14$
(maximum of $x - x^2$ at $x = \frac12$).

$\overline B(0, 1) = \{f : \sup\abs f \leq 1\}$: the continuous
functions with values in $\intcc{-1}{1}$.

$\{f : f(0) = 0\}$ is the preimage of $\{0\}$ under the
\emph{evaluation} $f \mapsto f(0)$, which is $1$-Lipschitz
($\abs{f(0) - g(0)} \leq d_\infty(f,g)$), hence continuous: the set
is closed (\cref{thm:b2:metric:globalcontinuity}). Likewise $\{f :
f(0) > 0\}$ is the preimage of the open $\intoo{0}{+\infty}$: open.
\end{solution}

\begin{solution}{exo:b2:metric:4}
\emph{Complete:} a Cauchy sequence with $\varepsilon = \frac12$ is
eventually constant, hence convergent.

\emph{Compact iff finite:} if $X$ is finite, any sequence takes some
value infinitely often (constant subsequence). If $X$ is infinite,
a sequence of pairwise distinct points has all mutual distances $1$:
no Cauchy subsequence, so no convergent one.

\emph{Connected subsets:} the singletons (and $\emptyset$). Any $A$
with two points $x \neq y$ splits as $\{x\} \cup (A
\setminus\{x\})$, both open in $A$ (every subset of a discrete
space is open --- balls of radius $\frac12$ are singletons).
\end{solution}

\begin{solution}{exo:b2:metric:5}
The function $g(x) = d(x, f(x))$ is continuous on the compact $X$
($\abs{g(x) - g(y)} \leq 2d(x,y)$ by two triangle inequalities), so
it attains its minimum at some $a$
(\cref{thm:b2:metric:compactprops}). If $f(a) \neq a$:
\[
g\bigl(f(a)\bigr) = d\bigl(f(a), f(f(a))\bigr) < d(a, f(a)) = g(a),
\]
contradicting minimality. So $f(a) = a$; uniqueness as usual (two
fixed points $a \neq b$ give $d(a,b) = d(f(a), f(b)) < d(a,b)$).

Non-compact example: $f(x) = x + \frac1x$ on $\intco{1}{+\infty}$:
$\abs{f(x) - f(y)} = \abs{x - y}\,\abs{1 - \frac{1}{xy}} < \abs{x -
y}$ for $x \neq y$ (as $xy > 1$), yet $f(x) > x$ everywhere.
\end{solution}

\begin{solution}{exo:b2:metric:6}
Pick $x_n \in K_n$. All terms from rank $n$ on lie in $K_n$; in
particular the whole sequence lies in the compact $K_0$: extract
$x_{\varphi(k)} \to \ell$. For each fixed $n$, the terms
$x_{\varphi(k)}$ with $\varphi(k) \geq n$ lie in the \emph{closed}
$K_n$, so the limit $\ell \in K_n$. Hence $\ell \in \bigcap K_n$.

Closed counterexample: $F_n = \intco{n}{+\infty}$ in $\R$: nested,
closed, nonempty, empty intersection.
\end{solution}

\begin{solution}{exo:b2:metric:7}
Continuity of $f^{-1}$ means: images $f(F)$ of closed sets $F
\subseteq K$ are closed (preimages under $f^{-1}$ are images under
$f$). A closed $F$ in the compact $K$ is compact
(\cref{thm:b2:metric:compactprops}~(1)); its continuous image
$f(F)$ is compact, hence closed. So $f^{-1}$ is continuous: $f$ is
a homeomorphism.

Counterexample: $\gamma(t) = (\cos t, \sin t)$ from
$\intco{0}{2\pi}$ (not compact) onto the unit circle is a
continuous bijection, but $\gamma^{-1}$ is discontinuous at $(1,0)$:
points on the circle just below the axis have parameters near
$2\pi$, not near $0$.
\end{solution}

\begin{solution}{exo:b2:metric:8}
$C = \bigcap_n C_n$ where each $C_n$ (union of $2^n$ closed
intervals of length $3^{-n}$) is closed: $C$ is closed and bounded
in $\R$, hence compact (\cref{thm:b2:metric:compactprops}~(2)).

Empty interior: $C$ contains no interval of length $> 3^{-n}$ (it
sits inside $C_n$, whose components have that length), for every
$n$.

Cardinality: the points of $C$ are exactly the reals $\sum_{n\geq1}
a_n 3^{-n}$ with digits $a_n \in \{0, 2\}$ (at each stage, the
deleted middle third removes the digit $1$); the map $(a_n) \mapsto
\sum a_n 3^{-n}$ is a bijection from $\{0,2\}^{\N^*}$ onto $C$
(injectivity as in \cref{exo:b2:structures:3}). So $C$ is equipotent
to $\{0,1\}^{\N}$: uncountable, though of ``length zero''.
\end{solution}

\begin{solution}{exo:b2:metric:9}
Suppose $a \notin f(X)$. The image $f(X)$ is compact (continuous
image), hence closed; so
\[
\varepsilon = d\bigl(a, f(X)\bigr)
= \inf_{y \in f(X)} d(a, y) > 0
\]
(the infimum of a continuous function on a compact set is attained;
if it were $0$, $a$ would be adherent to the closed $f(X)$, hence in
it).

Consider the orbit $x_n = f^n(a)$ ($x_0 = a$). For $p < q$:
\[
d(x_p, x_q) = d\bigl(f^p(a), f^p(f^{q-p}(a))\bigr)
= d\bigl(a, f^{q-p}(a)\bigr) \geq \varepsilon,
\]
(isometry iterated $p$ times; and $f^{q-p}(a) \in f(X)$ since $q - p
\geq 1$). A sequence with mutual distances $\geq \varepsilon$ has no
convergent subsequence --- contradicting compactness. Hence $f(X) =
X$.
\end{solution}

\begin{solution}{exo:b2:metric:10}
Let $A, B \in GL_n(\C)$ and $p(z) = \det\bigl((1-z)A + zB\bigr)$: a
polynomial in $z$ (each entry is affine in $z$; the determinant is a
polynomial in the entries). $p(0) = \det A \neq 0$: $p$ is not
identically zero, so it has finitely many roots $z_1, \dots, z_m$
(none equal to $0$ or $1$: $p(1) = \det B \neq 0$). The plane $\C$
minus finitely many points is path-connected: a path from $0$ to $1$
avoiding the $z_i$ exists (take a broken line through a point far
from all roots, or a circular arc; only finitely many obstacles).
Along such a path $\gamma$, $t \mapsto (1 - \gamma(t))A +
\gamma(t)B$ is a continuous path \emph{inside} $GL_n(\C)$ from $A$
to $B$ (the determinant never vanishes on it). Hence $GL_n(\C)$ is
path-connected --- unlike its real cousin
(\cref{ex:b2:metric:glnr}): the complex plane has room to walk
around obstacles.
\end{solution}
